\section{Discussion}
\label{sec:discussion}

\paragraph{What the measurements establish.} In this regime (a single shared tied output projection, ${\leq}350$M, undertrained) the standard cosine
conflict diagnostic misreads the CE/MTP interaction. On the rows that carry a
target the per-row gradients are aligned (active-row median cosine ${\approx}0.53$,
${\approx}3$--$4\%$ opposed; the full-vocabulary median is $+1$ only because
${\approx}92\%$ of rows are scalar multiples by construction), the aggregate
is $\approx 0$ or negative, and the two decouple during training. Instrumenting the per-row norms shows the aggregate is set by a high-norm opposed minority ($\approx0.3\%$ of all rows, carrying $60$--$69\%$ of the mass), not by disjoint support. The losses load the same rows (norm-profile cosine $0.98$) but oppose on the few rows that dominate the weighted sum. Crucially, the next-token degradation is large and
unambiguous, yet neutralizing that directional opposition does not remove it. The surgery
baselines in Appendix~\ref{app:surgery} recover nothing, which is consistent
with the opposed minority not being the fixable lever. Under the soft-label reading
this is expected for a sharper reason than opposition that carries the mass yet is
causally inert: the objective's \emph{minimizer has moved} to the mixture $q^{*}$,
and reprojecting gradients does not move it back. For symmetric PCGrad the argument
is exact: at a stationary point of the summed loss $g_{\text{CE}} = -g_{\text{MTP}}$,
both project to zero, and the mixture optimum stays stationary. Our asymmetric
head-only Gram--Schmidt is not a fixed point of those dynamics in the same way, so
there the empirical null is the evidence, and the natural reading is that a head-only
correction cannot undo a shifted optimum while the trunk keeps receiving the full
mixture gradient and shaping representations toward $q^{*}$. Capacity separation helps precisely because the auxiliary no longer forces one set of logits to serve every horizon,
though in our runs it reduces rather than eliminates the gap.

\paragraph{The stop-gradient result cuts the other way.}
One intervention appears to contradict the shifted-optimum account, and we state it
plainly rather than leave it to a reader. If the head's blended target were the
whole story, then $\mathrm{sg}(W)h$ should help: the head is freed to minimize pure
next-token cross-entropy given the representation it is handed. Instead
stop-gradient is the worst of the auxiliary conditions in
Table~\ref{tab:conditions} ($4.493$, $\Delta{+}0.508$ against shared MTP's
$\Delta{+}0.392$); only the gradient-surgery and optimizer-coverage variants of
Appendix~\ref{app:surgery} sit higher. The reason we cannot read this as a refutation is
that the probe moves two quantities at once. It removes the head's mixture target,
which the account predicts should help, and it also stops the head from co-adapting
to a trunk that is still being shaped by the full mixture gradient, which should
hurt. The observed sign says the second effect is the larger one; it does not say
the first is absent. Separating them needs an intervention that removes the mixture
target from the head \emph{without} decoupling head and trunk, which is what the
separate-head variant in \emph{Future work} does. Until then we treat the
stop-gradient number as evidence against the naive head-detach fix and as
uninformative about the mixture account itself.

Beyond the diagnostic result, an exact rewrite shows the shared-head objective is
soft-label cross-entropy toward a next-and-future-token mixture
(Section~\ref{sec:softlabel}), and a zero-parameter estimate of the implied KL cost,
computed entirely from the CE-only model's own distributions, matches the observed
degradation within run noise at the lower auxiliary weights while the full-mixture
variant reproduces ${\approx}88$--$89\%$ of the matched within-sweep gap at full weight, leaving a
${\approx}0.04$-nat residual attributable to the estimator's remaining downward
proxy bias and/or a small genuine non-mixture component
(Section~\ref{sec:mixturetest}). The anchored weight sweep shows the degradation
scales with the auxiliary weight as this account predicts, as a weight-proportional
interference account would also predict; the two shapes are indistinguishable at our
noise level.

\paragraph{Two failure modes, one instrument.} The aggregate cosine misleads in
two opposite directions depending on the optimizer, and the paper's central claim
is that both are failures of the \emph{same} diagnostic. Under Muon the aggregate
is ${-}0.08$ to ${-}0.19$ across measurement batches, close enough to zero at its shallowest reading (${-}0.08$) that a conflict diagnostic reads
``no interference,'' a \emph{false negative}: the next-token loss degrades by
$0.39$ nats regardless. Under AdamW the aggregate is ${-}0.28$ (${-}0.30$ to
${-}0.34$ over the norm-committing seeds), negative enough to trip a conflict flag,
and the surgery runs are \emph{consistent with} a false positive: projecting out
exactly this opposition recovers nothing (Appendix~\ref{app:surgery}). Those runs are
single-seed, with a drift envelope ($0.02$--$0.04$ nats) that overlaps the effects
they would need to resolve, so this arm is suggestive rather than established; the
matched-seed replication that would settle it is future item~3 below. A practitioner reading only the head's aggregate would conclude
``safe to add MTP'' under Muon and ``apply surgery'' under AdamW, and both
conclusions are wrong, the first demonstrably and the second on suggestive
evidence. The false negative is specific to the head: as
Section~\ref{sec:general} notes, an all-parameter aggregate would very
likely read clearly negative under both optimizers, so the misread is of
\emph{where} interference lives rather than of whether any is present. The instrument
does not merely lose power; it points the wrong way, and
which way depends on an optimizer choice unrelated to whether interference is
present. This is why the per-row and per-norm decomposition, not a scalar
threshold on the aggregate, is what the regime requires.

\section{Limitations}
\label{sec:limitations}
We state the boundaries of the result under four headings: what the scope of the
claim is, what is not established about mechanism, what biases the estimator
carries, and where the statistics are thin.

\paragraph{Scope.} Every measurement is on a single shared tied output projection
at ${\leq}350$M parameters, in an undertrained, small-batch regime where MTP is
already documented to underperform \citep{gloeckle2024mtp,zuhri2025top}. We do not show that MTP
\emph{as practiced} (with separate per-horizon heads) degrades next-token
modeling. The reason separate heads should help is not that they stop the losses
from sharing rows: the standard designs share one unembedding matrix, so per-row
gradient mixing persists there. It is that per-horizon logits remove the
constraint that one distribution serve every horizon, which is precisely the
optimum shift of Equation~\ref{eq:softlabel}. That makes the separate-head
experiment a clean test of our own account rather than a mere scaling check, since
it holds shared rows fixed and varies only the shared-label ingredient. The
outcome measure is also narrow: we record next-token validation loss, not downstream
task performance, so ``degrades'' should be read as that quantity and no more.

\paragraph{Mechanism.} We do not establish causation between the norm-weighted opposition and the loss increase. We show they co-occur, and that interventions behave as a shifted-optimum account predicts: surgery that removes the opposition does not remove the degradation, and separating capacity helps.
A related gap is one of timing rather than of inference. The per-row and per-norm
measurements are taken early in training, at step $1{,}000$, while the loss
comparisons are at $30{,}517$ steps; late-training persistence of the alignment is
measured directly only at 350M, on a single seed. The aligned-gradient state and
the degraded-loss state are therefore never exhibited in the same measurement at
124M, and the connection between them rests on that persistence rather than on a
co-located observation.

\paragraph{Estimator.} The zero-parameter prediction sits below the observed gap
by $0.035$--$0.046$ nats: within its stated biases, but also within what a small
genuine interference residual would produce. We treat the mixture account as the
leading candidate, not as confirmed, and we name the residual rather than assign
it. The estimate is also truncated rather than fully marginalized, and truncation
biases it upward relative to full marginalization, so a fuller estimate would
lower the prediction and \emph{widen} the residual rather than close it.

\paragraph{Which gradient we measure.} The head is weight-tied, so the gradient our
instrument reads sums the output-projection path and the embedding path
(Section~\ref{sec:instrument}). The embedding path deposits only on rows that already
carry a target, so it lands exactly on the active rows every statistic here is
computed over, and we do not bound its share. Two measurements would close this, and neither needs new
training. One is to recompute the projection path alone, taking the gradient with
respect to the logits and forming $\sum_t \delta_i(t) h_t$ directly. The other is to
commit per-row norms from an untied run, where the head tensor receives the
projection path only. The existing untied run cannot serve, because it committed
full-vocabulary fractions rather than per-row arrays.

\paragraph{Statistics and instrument caveats.} Seed counts are small and uneven: $n{=}5$ for CE-only and the standard $G_{\text{nce}}$ auxiliary, $n{=}3$ for the other four conditions (Table~\ref{tab:conditions}). The auxiliary-weight sweep of Section~\ref{sec:mixturetest} is a single seed throughout, so the strongest quantitative claim in the paper carries no across-seed spread of its own. The untied-head check is one run at $500$ steps
against $30{,}517$ elsewhere. Finally, the full-vocabulary form of the CE-vs-L1
control is threshold-sensitive and we do not rest calibration on it. On the
committed diagnostic run the fraction of rows with $|\cos|$ above $0.1$, $0.2$,
$0.3$ and $0.5$ is $74.9\%$, $35.2\%$, $0.42\%$ and $0$: nearly $35\%$ of
rows fall between the $0.2$ and $0.3$ bins, so the above-$0.3$ figure is
sensitive to exactly where the cut is placed, and it reads $36.2\%$ on the sweep
snapshot. That snapshot committed only these binned fractions, not the underlying per-row cosines. We therefore cannot say whether the two runs differ by a small shift near this quantile or more substantially, and we do not claim the discrepancy is explained. The shape of that distribution is itself unexplained, and we flag it rather than
account for it. Two comparisons bound the problem. Against isotropic vectors in
$d{=}768$ the cosine standard deviation is $0.036$, so almost no rows should clear
$0.1$, yet $74.9\%$ do: the rows are far more aligned than chance. But the profile is
also \emph{sharper} than any low-dimensional-span model predicts. Suppose both gradients were random within a shared
$k$-dimensional subspace. Any $k$ that reproduces the
above-$0.1$ fraction (within $5$ points of the observed $74.9\%$) forces the ratio of
the above-$0.2$ to the above-$0.3$ fraction to at most ${\approx}2$. The observed
ratio is $84$. Larger $k$ can produce a sharp ratio, but only by putting far too
little mass above $0.1$. The
best-fitting single $k$ is $23$, and it still overpredicts the above-$0.3$ fraction by
a factor of ${\approx}37$. So neither a chance model nor a shared-span model
reproduces the cliff, and we do not have per-row cosines from this run to diagnose it.
The control's role is to show that an \emph{unrelated} loss pair yields a near-zero
aggregate ($0.002$), not to quantify a per-row distribution. The unexplained shape
therefore does not bear on the calibration claim. It remains a loose end: committing
per-row cosines for the control is the measurement that would close it. What both runs agree on is the property the control
exists to establish: the aggregate cosine is near zero either way ($+0.002$ and
$-0.011$). This is why the active-row control of Section~\ref{sec:calib} is the
one we report.

\paragraph{Future work.} Five directions would broaden the scope from ``the diagnostic misreads this
regime'' toward a causal, general statement.
\begin{enumerate}
\item[(1)] A \textbf{separate-head MTP} variant, to localize whether the
      degradation is intrinsic to head-sharing.
\item[(2)] A \textbf{co-located measurement} at $30{,}517$ steps plus a 124M
      CE-only per-row baseline, to replace the timing assumption
      noted under \emph{Mechanism} above with a direct observation, and $n{=}3$ at
      350M to turn the single-seed scale observation into a measured trend.
\item[(3)] \textbf{Matched-seed, single-pass gradient surgery} at $n{=}3$, to
      remove the numerical-drift caveat.
\item[(4)] A \textbf{tighter KL estimate} (full marginalization or larger $k$,
      and better-trained proxy models), to bound the $0.035$--$0.046$-nat
      residual from the prediction side rather than the measurement side. The
      highest-value single piece is \textbf{estimating $P_3$ directly} instead of
      substituting $P_2$. That is the one bracket arm whose direction we have not
      bounded. Our own offset-3 cross-entropies suggest the substitution biases the
      prediction downward, so removing it could shrink the residual rather than
      merely bound it.
\item[(5)] \textbf{Committing per-row cosines for the control} at both
      checkpoints, which is what would settle the threshold-sensitivity
      question raised under \emph{Statistics and instrument caveats} above. Both snapshots are preserved,
      but neither stores the per-row array.
\end{enumerate}

\section{Conclusion}
\label{sec:conclusion}
Cosine similarity between task gradients is the default instrument for detecting
interference in multi-task and auxiliary-loss training. We showed that at a language model's shared, tied output projection this instrument is misleading, in three respects. The per-row CE and MTP gradients are aligned in direction on the rows that carry a target (active-row median cosine ${\approx}0.53$; the full-vocabulary median is $+1$ because most rows are scalar multiples by construction). The aggregate cosine is driven negative by a high-norm opposed minority ($\approx0.3\%$ of rows carrying most of the mass), not by a broad directional conflict. And the two metrics decouple during training.
Adding the shared auxiliary degrades next-token loss substantially, and the
surgery baselines we ran, single-seed and reported in
Appendix~\ref{app:surgery}, do not repair it. The aggregate is
set by per-row gradient norms, not by the per-row directions the cosine averages,
a distinction that explains why conflict-targeted surgery does not help here and
why capacity separation does. An exact rewrite of the objective supplies the
complementary positive account: the shared head optimizes cross-entropy against a
next-and-future-token mixture, so its optimum has moved, and a zero-parameter estimate of the
resulting KL cost reproduces most of the measured degradation, undershooting it by
$0.035$--$0.046$ nats at full auxiliary weight. On this reading the degradation is a predictable consequence of the objective rather than a pathology of optimization. Gradient surgery is then inert: exactly so for symmetric projection, and empirically so for the head-only variants we test, where the trunk continues to receive the full mixture gradient. We deliver the result
on a fully reproducible testbed with a committed statistics script, and we are
explicit about the shared-head, small-scale scope and the experiments that would
extend it.
